\section{Experimental Setup}
\label{sec:setup}

\subsection{Data}

Questions were drawn from SPIQA test-A~\cite{pramanick2024spiqa} (118 papers, 666
questions) with a fixed seed. Whole papers were assigned to the splits, so that no
paper contributes to both. A development split of 20 questions from 5 papers was used
for every design decision. A test split of 150 questions from 31 papers (71 with a
table and 79 with a figure as reference, at most 8 per paper) was used once. All
settings were frozen before the test run, and the evaluation script refuses to run on
the test split if any setting differs. A multi-hop subset of \nMultihopN{} test
questions was selected by a fixed rule over the question text: it holds the questions
that ask for a reason or an implication (``why'', ``suggest'', ``possible
explanation'') of a result shown in a table or figure.

\subsection{Systems Compared}

The same model, tools, prompts and grounding are used by every system, and each system
removes one part of the design. \emph{Full} is the complete agent. \emph{Single-pass}
has no planner, loops or gates: every tool is called once with the question, and the
writer answers from the result, so that the effect of agency is isolated.
\emph{Text-only} is the full graph with the text tool alone, with tables and figures
available only as extracted text, so that the effect of seeing them is isolated.
\emph{Strict G1} is the first frozen design, in which the agent abstains whenever G1
is not convinced (Section~\ref{sec:best-effort}). \emph{Full (Gemini)} is the full graph
with Gemini~3.5 Flash-Lite (free tier) behind every model call, so that the effect of
the model is separated from that of the architecture.

\subsection{Implementation}

The language model is Qwen2.5-VL-7B-Instruct~\cite{bai2025qwen25vl}, quantised to
4~bits and served by llama.cpp. All experiments were run offline on one
PC with an Intel Arc B580 GPU (12~GB) and 15.6~GB of RAM. Every model reply is
constrained to a JSON schema, generated at temperature 0 with a fixed seed, and cached
under a hash of the request, so that every answer can be reproduced exactly.

\subsection{Outcomes and Judge}

One of three outcomes is assigned to each question. It is \emph{correct} or
\emph{wrong} when at least one claim reaches the reader, and \emph{abstained} when the
system reports that it could not confirm an answer or when every claim is withheld by
G2. Answers were graded against SPIQA's gold answer by the local model, which was told
to accept an answer that gives the main point of the gold answer and contradicts
nothing in it, and to write its reason before its verdict. To check the judge, 30
answers (6 per system, shuffled) were graded by a human annotator who did not know the
system or the judge's verdict.

\subsection{Metrics}

Besides the outcome rates and the precision when answering, four scores were computed.
Retrieval faithfulness (RF) is the share of the distinct tables and figures in the
buffer that are the reference, and the hit rate is the share of questions whose
reference reached the buffer. The agent efficiency score (AES) is the number of hops
needed (one, or two on the multi-hop subset) divided by the number of search calls,
capped at one. The cross-modal grounding score (CGS) is the best intersection over union
(IoU) between a proof box and a hand-drawn answer box, with zero for a question
without a box. Since IoU was found to penalise correct proofs drawn at another size
(Section~\ref{sec:rq3}), the pointing game~\cite{zhang2018pointing} is also reported: a
proof counts as a hit when the centre of its box lies inside an answer box. This measure
was added after the IoU results were seen. Finally, the share of published claims that
cite the reference table or figure is reported as reference attribution. All rates are
given with 95\% bootstrap intervals over questions (10{,}000
resamples), and differences are paired by question.

\subsection{Annotation}

For the grounding scores, 50 test questions (25 tables and 25 figures) were drawn at
random before any output was inspected. On each page, the part of the table or figure
holding the gold answer was boxed by a co-author with an offline annotation tool and
written guidelines. In two quality-control rounds, 16 boxes that covered a whole table
or figure although the question named specific values were narrowed; they were
selected by a fixed rule, without reference to any system's boxes.
